\documentclass[10pt,a4paper]{amsart}
\usepackage[margin=19mm]{geometry}
\usepackage{amsmath,amssymb,mathtools}
\usepackage[scr=rsfs,cal=euler]{mathalfa}
\usepackage{xfrac}
\usepackage[unicode,hidelinks,bookmarks=false]{hyperref}
\newcommand{\ie}{i.e.\ }
\newcommand{\cf}{cf.\ }
\newcommand{\resp}{respectively\ }
\newcommand{\Subsection}{\subsection}
\providecommand{\nobreakdash}{\nobreak\hbox{-}}
\setlength{\emergencystretch}{2em}
\multlinegap0pt
\usepackage{amscd,latexsym,mathrsfs,bm,stackrel}
\usepackage{xy}
\usepackage{ulem}
\usepackage{tikz}
\usepackage{latexsym}
\usepackage{epsfig}
\newtheorem{theorem}{Theorem}[section]
\newtheorem{corollary}[theorem]{Corollary}
\newtheorem{lemma}[theorem]{Lemma}
\newtheorem{proposition}[theorem]{Proposition}
\newtheorem{conjecture}[theorem]{Conjecture}
\newtheorem{remark}[theorem]{\sc Remark}
\newtheorem{remarks}[theorem]{\sc Remarks}
\newtheorem{definition}[theorem]{Definition}
\newtheorem{example}[theorem]{\sc Example}
\newtheorem{question}[theorem]{\sc Question}
\newtheorem{note}[theorem]{\sc Note}
\newtheorem{comments}[theorem]{\sc Comments}
\numberwithin{equation}{section}  % numbers as (1.2) instead of (1) etc
\def\be{\begin{equation}}
\def\ee{\end{equation}}
\def\bt{\begin{theorem}}
\def\et{\end{theorem}}
\newcommand{\lra}{\longrightarrow}
\def\bc{\begin{corollary}}
\def\ec{\end{corollary}}
\def\br{\begin{remark}}
\def\er{\end{remark}}
\def\bex{\begin{example}}
\def\eex{\end{example}}
\def\bp{\begin{proposition}}
\def\ep{\end{proposition}}
\usepackage{amsbsy}
\usepackage{alltt,dsfont}
\usepackage{soul}
\usepackage{enumitem}
\usepackage{setspace}
\usepackage{tikz-cd}
\usepackage{textcomp}
\usepackage{colonequals}
\usepackage{mathtools}
\usepackage{bm}
\usepackage{mathptmx}
\newcommand{\C}{\mathbb{C}}
\newcommand{\Z}{\mathbb{Z}}
\newcommand{\Q}{\mathbb{Q}}
\newcommand{\R}{\mathbb{R}}
\newcommand{\N}{\mathbb{N}}
\newcommand{\K}{\mathbb{K}}
\renewcommand{\P}{\mathbb{P}}
\newcommand{\D}{\mathbb{D}}
\newcommand{\cA}{\mathcal{A}}
\newcommand{\cB}{\mathcal{B}}
\newcommand{\cC}{\mathcal{C}}
\newcommand{\cD}{\mathcal{D}}
\newcommand{\cE}{\mathcal{E}}
\newcommand{\cF}{\mathcal{F}}
\newcommand{\cG}{\mathcal{G}}
\newcommand{\cH}{\mathcal{H}}
\newcommand{\cI}{\mathcal{I}}
\newcommand{\cJ}{\mathcal{J}}
\newcommand{\cK}{\mathcal{K}}
\newcommand{\cL}{\mathcal{L}}
\newcommand{\cM}{\mathcal{M}}
\newcommand{\cN}{\mathcal{N}}
\newcommand{\cO}{\mathcal{O}}
\newcommand{\cP}{\mathcal{P}}
\newcommand{\cQ}{\mathcal{Q}}
\newcommand{\cR}{\mathcal{R}}
\newcommand{\cS}{\mathcal{S}}
\newcommand{\cT}{\mathcal{T}}
\newcommand{\cU}{\mathcal{U}}
\newcommand{\cV}{\mathcal{V}}
\newcommand{\cW}{\mathcal{W}}
\newcommand{\cX}{\mathcal{X}}
\newcommand{\cY}{\mathcal{Y}}
\newcommand{\cZ}{\mathcal{Z}}
\newcommand{\bA}{\mathbb{A}}
\newcommand{\bB}{\mathbb{B}}
\newcommand{\bC}{\mathbb{C}}
\newcommand{\bD}{\mathbb{D}}
\newcommand{\bE}{\mathbb{E}}
\newcommand{\bF}{\mathbb{F}}
\newcommand{\bG}{\mathbb{G}}
\newcommand{\bH}{\mathbb{H}}
\newcommand{\bI}{\mathbb{I}}
\newcommand{\bJ}{\mathbb{J}}
\newcommand{\bK}{\mathbb{K}}
\newcommand{\bL}{\mathbb{L}}
\newcommand{\bM}{\mathbb{M}}
\newcommand{\bN}{\mathbb{N}}
\newcommand{\bO}{\mathbb{O}}
\newcommand{\bP}{\mathbb{P}}
\newcommand{\bQ}{\mathbb{Q}}
\newcommand{\bR}{\mathbb{R}}
\newcommand{\bS}{\mathbb{S}}
\newcommand{\bT}{\mathbb{T}}
\newcommand{\bU}{\mathbb{U}}
\newcommand{\bV}{\mathbb{V}}
\newcommand{\bW}{\mathbb{W}}
\newcommand{\bX}{\mathbb{X}}
\newcommand{\bY}{\mathbb{Y}}
\newcommand{\bZ}{\mathbb{Z}}
\newcommand{\fa}{\mathfrak{a}}
\newcommand{\fb}{\mathfrak{b}}
\newcommand{\fc}{\mathfrak{c}}
\newcommand{\fd}{\mathfrak{d}}
\newcommand{\fe}{\mathfrak{e}}
\newcommand{\ff}{\mathfrak{f}}
\newcommand{\fg}{\mathfrak{g}}
\newcommand{\fh}{\mathfrak{h}}
\newcommand{\fj}{\mathfrak{j}}
\newcommand{\fk}{\mathfrak{k}}
\newcommand{\fl}{\mathfrak{l}}
\newcommand{\fm}{\mathfrak{m}}
\newcommand{\fn}{\mathfrak{n}}
\newcommand{\fo}{\mathfrak{o}}
\newcommand{\fp}{\mathfrak{p}}
\newcommand{\fq}{\mathfrak{q}}
\newcommand{\fr}{\mathfrak{r}}
\newcommand{\fs}{\mathfrak{s}}
\newcommand{\ft}{\mathfrak{t}}
\newcommand{\fu}{\mathfrak{u}}
\newcommand{\fv}{\mathfrak{v}}
\newcommand{\fw}{\mathfrak{w}}
\newcommand{\fx}{\mathfrak{x}}
\newcommand{\fy}{\mathfrak{y}}
\newcommand{\fA}{\mathfrak{A}}
\newcommand{\fB}{\mathfrak{B}}
\newcommand{\fC}{\mathfrak{C}}
\newcommand{\fD}{\mathfrak{D}}
\newcommand{\fE}{\mathfrak{E}}
\newcommand{\fF}{\mathfrak{F}}
\newcommand{\fG}{\mathfrak{G}}
\newcommand{\fH}{\mathfrak{H}}
\newcommand{\fI}{\mathfrak{I}}
\newcommand{\fJ}{\mathfrak{J}}
\newcommand{\fK}{\mathfrak{K}}
\newcommand{\fL}{\mathfrak{L}}
\newcommand{\fM}{\mathfrak{M}}
\newcommand{\fN}{\mathfrak{N}}
\newcommand{\fO}{\mathfrak{O}}
\newcommand{\fP}{\mathfrak{P}}
\newcommand{\fQ}{\mathfrak{Q}}
\newcommand{\fR}{\mathfrak{R}}
\newcommand{\fS}{\mathfrak{S}}
\newcommand{\fT}{\mathfrak{T}}
\newcommand{\fU}{\mathfrak{U}}
\newcommand{\fV}{\mathfrak{V}}
\newcommand{\fW}{\mathfrak{W}}
\newcommand{\fX}{\mathfrak{X}}
\newcommand{\fY}{\mathfrak{Y}}
\newcommand{\sA}{\mathscr{A}}
\newcommand{\sB}{\mathscr{B}}
\newcommand{\sC}{\mathscr{C}}
\newcommand{\sD}{\mathscr{D}}
\newcommand{\sE}{\mathscr{E}}
\newcommand{\sF}{\mathscr{F}}
\newcommand{\sG}{\mathscr{G}}
\newcommand{\sH}{\mathscr{H}}
\newcommand{\sI}{\mathscr{I}}
\newcommand{\sJ}{\mathscr{J}}
\newcommand{\sK}{\mathscr{K}}
\newcommand{\sL}{\mathscr{L}}
\newcommand{\sM}{\mathscr{M}}
\newcommand{\sN}{\mathscr{N}}
\newcommand{\sO}{\mathscr{O}}
\newcommand{\sP}{\mathscr{P}}
\newcommand{\sQ}{\mathscr{Q}}
\newcommand{\sR}{\mathscr{R}}
\newcommand{\sS}{\mathscr{S}}
\newcommand{\sT}{\mathscr{T}}
\newcommand{\sU}{\mathscr{U}}
\newcommand{\sV}{\mathscr{V}}
\newcommand{\sW}{\mathscr{W}}
\newcommand{\sX}{\mathscr{X}}
\newcommand{\sY}{\mathscr{Y}}
\newcommand{\of}{\overline{f}}
\newcommand{\oX}{\overline{X}}
\newcommand{\oF}{\overline{F}}
\newcommand{\ok}{\overline{k}}
\newcommand{\oi}{\overline{i}}
\newcommand{\oY}{\overline{Y}}
\newcommand{\oE}{\overline{E}}
\newcommand{\og}{\overline{g}}
\newcommand{\oG}{\overline{G}}
\newcommand{\oH}{\overline{H}}
\newcommand{\wti}{\widetilde}
\renewcommand{\d}{\partial}
\newcommand{\Gr}{\mathrm{Gr}}
\newcommand{\gr}{\mathrm{Gr}}
\newcommand{\sd}{\mathrm{sd}}
\newcommand{\Pol}{\mathrm{Pol}}
\newcommand{\IC}{\mathrm{IC}}
\newcommand{\DR}{\mathrm{DR}}
\newcommand{\Sing}{\mathrm{Sing}}
\newcommand{\Sym}{\mathrm{Sym}}
\newcommand{\MHM}{\mathrm{MHM}}
\newcommand{\MH}{\mathrm{MH}}
\newcommand{\Char}{\mathrm{Char}}
\newcommand{\bl}{\mathrm{bl}}
\newcommand{\Alb}{\mathrm{Alb}}
\DeclareMathOperator{\Ker}{Ker}
\DeclareMathOperator{\rk}{rank}
\DeclareMathOperator{\im}{Image}
\newcommand{\pr}{\mathrm{pr}}
\newcommand{\Sp}{\mathrm{Sp}}
\newcommand{\Hom}{\mathrm{Hom}}
\DeclareMathOperator{\supp}{supp}
\newcommand{\mhm}{\mathrm{MHM}}
\newcommand{\td}{\mathrm{td}}
\newcommand{\MHS}{\mathrm{mHs}}
\newcommand{\rat}{{\it {rat}}}
\newcounter{tmp}
\begin{document}
\title[On the topology of fibers of complex polynomial maps]{On the topology of fibers of complex polynomial maps}
\author{Laurentiu Maxim, John Messina}
\date{}
\maketitle
\noindent\emph{Source and licence.} On the topology of fibers of complex polynomial maps. Version arXiv:2608.04320v2 (CC BY 4.0). This material is distributed under the Creative Commons Attribution 4.0 International licence, http://creativecommons.org/licenses/by/4.0/, as recorded in the arXiv OAI licence metadata for this exact version and shown on the abstract page https://arxiv.org/abs/2608.04320v2. Full rights notice: licenses/arxiv-2608.04320-CC-BY-4.0.txt.\par\smallskip
\noindent\textit{Editorial scope.} Counted unit: the original \texttt{theorem} environment \texttt{lsc} at source lines 1253--1260 together with its complete proof at lines 1262--1322; the delivered document keeps source lines 1239--1323 so that every referenced label and every citation resolves inside it. Native editable source is the paper's own concatenated TeX; the AMS wrapper is editorial formatting only. No publisher class, upstream script or unrelated artwork is redistributed.
\section{Top Betti number of general fiber. Semi-continuity properties}\label{top}
Recall from \cite{T2} that a complex polynomial $f\colon\C^{n+1} \to \C$ is called {\it generic-at-infinity} (or, of $\mathcal{G}$-type) if, for any $c\in \C$, the compactified fiber $\oF_c \subset \C P^{n+1}$ is transversal to the hyperplane at infinity $H_\infty$. It can be shown that fibers of a $\cG$-type polynomial may have at most isolated singularities. %One important example is that of the Fermat polynomial $f(x_1, \ldots, x_{n+1})=\sum_{i=1}^{n+1} x_i^d.$

Degree $d$ polynomials of $\cG$-type are completely characterized by the fact that the top Betti number $b_n$ of their general fiber equals $(d-1)^{n+1}$, e.g., see  \cite[Prop.4.1.5]{T2}. The lower semi-continuity property for the top Betti number of general fibers (cf. \cite[Prop.2.1]{ST1}, \cite[Prop.4.2.1]{T2}) then yields the following upper bound result.

\begin{proposition}\label{pr61}
Let $f:\C^{n+1} \to \C$ be a degree $d$ complex polynomial with bifurcation set $B$. For any $c \in \C \setminus B$, one has
\be\label{tb}
b_n(F_c)\leq (d-1)^{n+1}.
\ee
\end{proposition}

In this section, we give a sheaf theoretic proof of the above-mentioned lower semicontinuity property for the top Betti number of general fibers, which has the advantage that it is also amenable to Hodge-theoretic considerations. (The homological version of this result is proved in \cite{ST1,T2} by homotopy-theoretic methods.)


\begin{theorem}\label{lsc}
%Let $f,g: \C^{n+1} \to \C$ be degree $d$ polynomials, with $P(x,s):=f(x)-sg(x):\C^{n+1} \times \C \to \C$ a one-parameter deformation of $f=P(-,0)$ by $g$. 
Let $P(x,s):\C^{n+1} \times \C \to \C$ be a one-parameter family of degree $d$ polynomials, %with $f=f_0=P(-,0)$.
and denote by $F^s$ the general fiber of $f_s:=P(-,s)$, $s \in \C$. Then there is an injective homomorphism 
\be\label{injl} H^n_c(F^0;\Q)  \rightarrowtail H^n_c(F^s;\Q) \ee
for $s\neq 0$ close enough to $0$. In particular, the top Betti number $b_n$ of general fibers satisfies the lower semi-continuity property, namely,
\begin{center} $b_n(F^s) \geq b_n(F^0)$, for $s \neq 0$ close enough to $0$.\end{center} 
\end{theorem}

\begin{proof}
For $c \in \C$ general enough, we may assume that $f_s^{-1}(c)=F^s$ is the general fiber of $f_s$, for $s$ in a small enough neighborhood of $0$. Consider the variety
\[
\overline{Y}:=\{\left([x_0:x],s\right) \in  \bC P^{n+1} \times \bC \mid P^h(x_0,x,s)-cx_0^d=0  \},
\]
where $P^h$ is the homogenization of $P$ by the variable $x_0$, considering $s$ as a parameter. Let $\pi:\overline{Y} \to \C$ be the projection to the $s$-coordinate. Note that $\oF^s:=\pi^{-1}(s)$ is the projective closure in $\C P^{n+1}$ of the affine hypersurface $F^s=f_s^{-1}(c)$.

Set \[ Y = \overline{Y}\cap \{x_0=1\}=\{(x,s) \in  \bC^{n+1} \times \bC \mid P(x,s)=c  \},\] with the open affine inclusion $u: Y\hookrightarrow \overline{Y}$, and note that $F^s=\oF^s \cap Y$.

Let $\overline{i}_0: \overline{F}^0 \hookrightarrow \overline{Y}$ be the inclusion map, and consider the distinguished triangle 
\[ \overline{i}_0^{*}u_{!}\Q_Y \lra \psi_{\pi}u_{!}\Q_Y \lra \varphi_{\pi}u_{!}\Q_Y \overset{[1]}{\lra}\]
and the following part of the associated hypercohomology long exact sequence:
\be\label{les10} \cdots \to \bH^{n-1}(\overline{F}^0; \varphi_{\pi}u_{!}\Q_Y) \to \bH^n(\overline{F}^0; \overline{i}_0^{*}u_{!}\Q_Y) \to \bH^{n}(\overline{F}^0; \psi_{\pi}u_{!}\Q_Y) \to \cdots \ee

Let $i_0:F^0 \hookrightarrow Y$ and $u_0:F^0 \hookrightarrow \overline{F}^0$ be the inclusion maps. 
Using proper base change, we have
\begin{equation}\label{les11}
\bH^n(\overline{F}^0; \overline{i}_0^{*}u_{!}\Q_Y) \cong 
   \mathbb{H}^n(\overline{F}^0; {u_0}_!{i}_0^{*}\mathbb{Q}_{Y}) \cong  \mathbb{H}^n(\overline{F}^0; {u_0}_!\mathbb{Q}_{F^0})  \cong H_c^n(F^0;\mathbb{Q}).
\end{equation}

Similarly, let $\overline{i}_{s}: \overline{F}^{s}\hookrightarrow \overline{Y}$, $u_{s}: F^{s} \hookrightarrow \overline{F}^{s}$ and $i_{s}: {F}^{s}\hookrightarrow Y$ be the inclusions. Then, using \cite[Ex.10.4.20]{MS} and proper base change, for $s\neq 0$ sufficiently close to $0$ we have
\begin{equation}\label{les12}
\bH^{n}(\overline{F}^0; \psi_{\pi}u_{!}\Q_Y) \cong 
 \mathbb{H}^n(\overline{F}^{s}; \overline{i}_{s}^{*}u_{!}\mathbb{Q}_{Y})
 \cong \mathbb{H}^n(\overline{F}^{s}; {u_{s}}_!i_{s}^{*}\mathbb{Q}_{Y}) \cong H^n_c(F^{s};\mathbb{Q}).
\end{equation}

Therefore, by combining \eqref{les10}, \eqref{les11} and \eqref{les12}, we note that in order to get an injection $H^n_c(F^0;\Q)  \rightarrowtail H^n_c(F^s;\Q)$ as in the statement of the theorem, it suffices to show that  $\bH^{n-1}(\overline{F}^0; \varphi_{\pi}u_{!}\Q_Y) \cong 0$.


Note that
\begin{equation}\label{lab1}
\bH^{n-1}(\overline{F}^0; \varphi_{\pi}u_{!}\Q_Y) \cong 
\mathbb{H}^{-1}(\overline{F}^0;{}^p\varphi_{\pi}u_{!}\mathbb{Q}_{Y}[n+1])
\cong H^{-1}(R{\pi_0}_*{}^p\varphi_{\pi}u_{!}\mathbb{Q}_{Y}[n+1]),
\end{equation}
where ${}^p\varphi_{\pi}=\varphi_{\pi}[-1]$ is the perverse vanishing cycle functor,  and $\pi_0\colon \overline{F}^0\rightarrow \{0\}$ is the constant map.
 Now, consider the maps $\overline{Y}\xrightarrow{\pi}\mathbb{C}\xrightarrow{\textrm{id}}\mathbb{C},$ with induced maps $\overline{F}^0\xrightarrow{\pi_0} \{0\}\xrightarrow{\textrm{id}} \{0\}$ of zero sets. Since $\pi$ is proper, by proper base change for vanishing cycles (e.g.,  see \cite[Rem.4.3.7]{Sch}) we have
\begin{equation}
    \varphi_{\textrm{id}}\circ R\pi_{*}\simeq R{\pi_{0}}_*\circ \varphi_{\pi}
\end{equation}
and so
\begin{equation}\label{lab2}
H^{-1}(R{\pi_0}_*{}^p\varphi_{\pi}u_{!}\mathbb{Q}_{Y}[n+1])\cong H^{-1}({}^p\varphi_{\textrm{id}} R\pi_{*}u_{!}\mathbb{Q}_{Y}[n+1]) \cong H^{-1}({}^p\varphi_{\textrm{id}}R(\pi \circ u)_{!}\mathbb{Q}_Y[n+1]),
\end{equation}
where the last isomorphism uses the properness of $\pi$. 

Since $Y$ is a complex hypersurface of dimension $n+1,$ the complex $\mathbb{Q}_Y[n+1]$ is perverse on $Y$. Since $\pi \circ u : Y \to \C$ 
is an affine map, the functor $R(\pi \circ u)_{!}$ is left $t$-exact, and so $R(\pi \circ u)_{!}\mathbb{Q}_Y[n+1]\in {}^p D^{\geq 0}(\mathbb{C}).$ Then since ${}^p\varphi_{\textrm{id}}$ is $t$-exact, we get that ${}^p\varphi_{\textrm{id}} R(\pi \circ u)_{!}\mathbb{Q}_Y[n+1]\in {}^pD^{\geq 0}(\{0\})=D^{\geq 0}(\{0\}),$ which yields that 
\begin{equation}\label{lab3}
    H^{-1}({}^p\varphi_{\textrm{id}}R(\pi \circ u)_{!}\mathbb{Q}_Y[n+1])=0.
\end{equation}
Then by \eqref{lab1},  \eqref{lab2} and \eqref{lab3} we get that
\begin{equation}
    \bH^{n-1}(\overline{F}^0; \varphi_{\pi}u_{!}\Q_Y) \cong 0,
\end{equation}
as desired.

Finally, since both $F_0$ and $F_{s}$ are smooth of complex dimension $n$, we get from \eqref{injl} by Poincar\'e duality that \[b_n(F^0) \leq b_n(F^{s}).\]
\end{proof}


\begin{thebibliography}{99}
\bibitem[MS]{MS} Lauren\c{t}iu Maxim, J\"org Sch\"urmann, {\it Constructible sheaf complexes in complex geometry and applications}, Handbook of geometry and topology of singularities III, 679--791, Springer, Cham, 2022.
\bibitem[ST1]{ST1} Dirk Siersma, Mihai Tib\u{a}r, {\it Singularity exchange at the frontier of the space}, Real and complex singularities, 327--342, Trends Math., Birkh\"auser, Basel, 2007.
\bibitem[Sch]{Sch} J\"org Sch\"urmann, {\it Topology of singular spaces and constructible sheaves}, Mathematics Institute of the Polish Academy of Sciences. Mathematical Monographs (New Series), 63. Birkh\"auser Verlag, Basel, 2003.
\bibitem[T2]{T2} Mihai Tib\u{a}r, {\it Polynomials and vanishing cycles.} Cambridge Tracts in Mathematics, 170. Cambridge University Press, Cambridge, 2007.
\end{thebibliography}
\end{document}
